\documentclass[11pt]{article}
\usepackage[margin=1in]{geometry}
\usepackage{mathptmx}
\usepackage{courier}
\usepackage[T1]{fontenc}
\usepackage[utf8]{inputenc}
\usepackage{amsmath,amssymb,amsthm,mathtools}
% --- Robust definitions to work around mathptmx/hyperref issues ---
\let\oldhbar\hbar
\DeclareRobustCommand{\hbar}{\ensuremath{\hslash}}
\newcommand{\SRa}{S_{\mathrm{RA}}}
\newcommand{\Dkl}{D_{\mathrm{KL}}}
\newcommand{\TV}{\mathrm{TV}}
% --- End robust definitions ---

\usepackage{enumitem}
\usepackage{hyperref}
\usepackage{xcolor}
\usepackage{setspace}
\usepackage{titlesec}
\usepackage{booktabs}
\usepackage{tcolorbox}
\usepackage{array}
\hypersetup{colorlinks=true,linkcolor=blue,citecolor=blue,urlcolor=blue}

% Prevent math commands from breaking PDF metadata (bookmarks, abstract title, etc.)
\pdfstringdefDisableCommands{%
  \def\hbar{hbar}%
  \def\ell{l}%
  \def\Delta{Delta}%
  \def\Lambda{Lambda}%
  \def\Omega{Omega}%
  \def\rho{rho}%
  \def\sigma{sigma}%
  \def\mu{mu}%
  \def\Pi{Pi}%
  \def\alpha{alpha}%
  \def\beta{beta}%
  \def\gamma{gamma}%
  \def\tau{tau}%
  \def\theta{theta}%
  \def\kappa{kappa}%
  \def\lambda{lambda}%
  \def\nu{nu}%
  \def\pi{pi}%
  \def\phi{phi}%
  \def\psi{psi}%
  \def\chi{chi}%
  \def\xi{xi}%
  \def\eta{eta}%
  \def\zeta{zeta}%
  \def\varepsilon{epsilon}%
  \def\varphi{phi}%
  \def\mathrm#1{#1}%
  \def\mathbf#1{#1}%
  \def\mathcal#1{#1}%
  \def\mathbb#1{#1}%
  \def\text#1{#1}%
  \def\sqrt#1{sqrt(#1)}%
  \def\frac#1#2{#1/#2}%
  \def\cdot{*}%
  \def\times{x}%
  \def\leq{<=}%
  \def\geq{>=}%
  \def\approx{~}%
  \def\to{->}%
  \def\rightarrow{->}%
  \def\leftarrow{<-}%
  \def\infty{inf}%
  \def\sum{sum}%
  \def\int{int}%
  \def\prod{prod}%
  \def\partial{d}%
  \def\nabla{grad}%
  \def\propto{prop}%
  \def\ldots{...}%
  \def\cdots{...}%
  \def\dots{...}%
  \def\pm{+/-}%
  \def\SRa{S_RA}%
  \def\Pacc{P_acc}%
  \def\Dkl{D_KL}%
  \def\TV{TV}%
}

\setstretch{1.08}
\titleformat{\section}{\large\bfseries}{\thesection.}{0.6em}{}
\titleformat{\subsection}{\normalsize\bfseries}{\thesubsection.}{0.5em}{}
\newtheorem{axiom}{Axiom}
\newtheorem{definition}{Definition}
\newtheorem{proposition}{Proposition}
\newtheorem{remark}{Remark}
\newtheorem{conjecture}{Conjecture}
\newtheorem{theorem}{Theorem}
\newtheorem{lemma}{Lemma}
\newtheorem{corollary}{Corollary}
\newcommand{\RA}{\mathrm{RA}}
\newcommand{\Pacc}{P_{\mathrm{acc}}}

\begin{document}
\begin{center}
{\LARGE \textbf{Relational Actualism IV:\\[0.3em] Complexity, Life, and the Causal Firewall}}\\[1em]
{\large Joshua F. Sandeman}\\[0.25em]
{\normalsize Independent Researcher, Salem, Oregon}\\[0.15em]
{\small sansuikyo@gmail.com}\\[0.5em]
{\normalsize Draft --- April 14, 2026}
\end{center}

\begin{abstract}
This paper extends Relational Actualism (RA) into the domains of complexity, life, agency, and cognition. Its central aim is not to claim complete derivation of these domains, but to develop a common discrete framework in which they can be discussed using the same ontology as the rest of the RA suite: finite actual graphs, structured potentia, shielding, recursive closure, and the causal firewall. The paper contains three distinct kinds of claims, which are explicitly distinguished throughout: structural results concerning causal organization and boundary formation; empirical or observational predictions concerning life, biosignatures, and the causal firewall; and exploratory philosophical implications concerning agency, perception, and consciousness. The strongest contributions of the paper concern recursive closure, shielding, and the organizational conditions under which stable complexity can persist. Where consciousness is discussed, the paper aims at structural necessary conditions and exclusion results, not a derived sufficient-condition theory of phenomenal consciousness. The paper should therefore be read as the most exploratory member of the RA suite: a disciplined extension of the framework into complexity and life, rather than a claim of full closure on these topics.
\end{abstract}

\begin{tcolorbox}[title=Epistemic Status of This Paper,colback=gray!5,colframe=gray!60]
This is the most exploratory paper in the RA suite. Its claims fall into three categories:
\begin{itemize}[nosep]
    \item \textbf{Structural results} concerning causal organization, recursive closure, shielding, and boundary formation;
    \item \textbf{Empirical predictions} concerning the causal firewall, biosignatures, and life-related observational signatures;
    \item \textbf{Exploratory philosophical implications} concerning agency, perception, and consciousness.
\end{itemize}
Where consciousness is discussed, the paper aims at structural necessary conditions and exclusion results, not a derived sufficient-condition theory of phenomenal consciousness.
\end{tcolorbox}

\begin{tcolorbox}[title=Place of This Paper in the RA Suite,colback=gray!5,colframe=gray!60]
This paper is the most exploratory member of the Relational Actualism suite. It extends the framework into complexity, life, agency, and the causal firewall.
\begin{itemize}[nosep]
    \item \textbf{Assumed from Paper~I:} the ontology, axioms, finite growth rule, and actualization framework.
    \item \textbf{Used from Paper~II where relevant:} the motif, renewal, and interaction structures required for stable complex organization.
    \item \textbf{Used from Paper~III where relevant:} severance, shielding, boundary structure, and the macroscopic interpretation of causal organization.
    \item \textbf{This paper contributes:} the complexity hierarchy, the causal firewall programme, and exploratory applications to life, perception, agency, and consciousness.
\end{itemize}
This paper should be read as a disciplined extension of the RA framework rather than as a claim of theorem-level closure on all questions concerning life or mind.
\end{tcolorbox}

%% ═══════════════════════════════════════════════════════════
\section{Introduction: From Bosons to Brains}
%% ═══════════════════════════════════════════════════════════

\paragraph{How to read this paper.}
This paper is the most exploratory member of the RA suite. Its purpose is to extend the discrete framework of Relational Actualism into the domains of complexity, life, agency, and cognition, but it does not claim theorem-level closure across all of these topics. Three different claim types appear below and are explicitly distinguished: first, structural claims about causal organization, recursive closure, shielding, and boundary formation within the RA graph; second, empirical or observational predictions concerning life, biosignatures, and the causal firewall; and third, exploratory philosophical implications concerning agency, perception, and consciousness. In particular, where consciousness is discussed, the paper aims at structural necessary conditions and exclusion results, not a derived sufficient-condition theory of phenomenal consciousness.

The Relational Actualism programme began with the claim that irreversible actualization events are the primitive physical reality. Papers~I--III develop the consequences of this commitment, filtered through the BDG integers $(1,-1,9,-16,8)$, in the domains of quantum mechanics, the Standard Model, general relativity, and cosmology. The discrete/topological core of those papers is well-supported; several continuum-bridge steps are explicitly labeled DR (derived) conditional, PI (phenomenological interpolation), or CN (conjecture), as their epistemic-status tables document. The present paper asks: does the same machinery extend to biological complexity?

The answer is yes, and the mechanism is specific. The Erd\H{o}s--R\'enyi percolation transition at actualization density $\mu=1$ is the mathematical boundary at which a random graph acquires a giant connected component. This transition appears in five physically distinct contexts within RA, all governed by the same critical parameter:
\begin{enumerate}[nosep]
\item QCD confinement: the transition from free quarks to hadrons (Paper~II).
\item Galactic rotation curves: the density threshold between 4D and 2D graph geometry.
\item Fault-tolerant quantum computing: the error-correction threshold $p_{\mathrm{th}}=\alpha_{\mathrm{EM}}$.
\item The actualization filter: maximum selectivity $\Delta S^*$ at $\mu=1$ (Paper~I).
\item The origin of biological complexity: the Causal Firewall percolation threshold.
\end{enumerate}

These are not five independent coincidences. They are five proposed expressions of the same geometric fact: the Poisson-CSG graph at $\mu=1$ is at the edge of connectivity. The five contexts differ, however, in how fully each identification has been closed. Some rest on tight derivations from the BDG filter (context 4, the actualization threshold itself; context 1, QCD confinement via BDG path-weight enumeration); others rest on a continuum-bridge identification (context 2, galactic rotation curves, which inherits a Paper~III dimensional-reduction bridge); others are programme-level identifications whose physical content is still interpretive (context 3, the quantum-computing threshold, and context 5, the causal-firewall percolation identification developed in this paper). The table in \S\ref{sec:status} tracks these distinctions; the $\mu=1$ architecture should be read as a proposed unifying thread across RA, not as an equally-closed identification in all five domains.

The chain from fundamental physics to biological complexity is proposed here as a unified architecture rather than a closed series of theorems. Several links are genuinely derived at the discrete layer; others are identifications under the BDG-to-metric or actualization-to-macroscopic dictionaries, still in need of tighter closure. The present paper makes the architecture explicit while documenting which links sit at which epistemic layer.

%% ═══════════════════════════════════════════════════════════
\section{The Complexity Hierarchy}
\label{sec:hierarchy}
%% ═══════════════════════════════════════════════════════════

The same coarse-graining operation that produces the spacetime metric from the actualized graph (Paper~III), applied recursively, produces four tiers of physical organization. This is not a hierarchy of new laws imposed at each scale. It is the \emph{same} growth rule viewed through successively coarser lenses.

\textbf{Tier~1: Fundamental.} Individual actualization events. Governed by the BDG action directly. Particle physics, quantum mechanics, and the measurement problem (Paper~I).

\textbf{Tier~2: Mesoscopic.} Coarse-grained clusters of actualization events forming stable local patterns---atoms, molecules, crystals. The LLC at this scale gives chemical conservation laws. Stable renewal motifs (Paper~II) are Tier~2 structures.

\textbf{Tier~3: Macroscopic.} Large-scale patterns of Tier~2 structures. Spacetime geometry, thermodynamics, classical mechanics. General relativity is the Tier~3 effective description of the actualized graph (Paper~III). The Milne/EdS expansion regimes are Tier~3 dynamics.

\textbf{Tier~4: Complex.} Self-reinforcing, self-replicating patterns of Tiers~2 and~3. Biological organisms, ecosystems, minds. The defining property is \emph{recursive closure}: a Tier~4 pattern maintains its own boundary conditions through its own dynamics.

\begin{remark}
The hierarchy is not a ladder of emergence in the strong sense (new laws at each level). It is a ladder of coarse-graining in the weak sense (the same law, viewed at different resolutions). Each tier inherits the LLC, the BDG filter, and the $\mu=1$ threshold from the tier below. What changes is the effective degree of freedom---from individual vertices to atoms to spacetime patches to organisms.
\end{remark}

The critical transition is between Tier~3 and Tier~4. Tier~3 structures (rocks, stars, galaxies) are stable but not self-replicating. Tier~4 structures (cells, organisms) maintain their own causal boundary through active metabolism---they are \emph{causally self-sustaining} in a way that Tier~3 structures are not. The Causal Firewall is the physical boundary between Tier~3 and Tier~4.

\subsection{Recursive coarse-graining formalism}
\label{sec:coarse_grain}

The tier hierarchy is not vague: it is a recursive coarse-graining operation applied to the causal graph. Let $G_0$ be the actualized causal graph at the Planck scale (Tier~1). The coarse-graining operation $\Phi$ partitions $G_0$ into clusters, each of which becomes a single vertex in the coarse-grained graph $G_1 = \Phi(G_0)$. The causal edges of $G_1$ are the coarse-grained images of the causal edges of $G_0$, preserving the partial order.

At each step, the operation preserves:
\begin{itemize}[nosep]
\item The local ledger condition: if $G_0$ satisfies LLC at every vertex, so does $G_1$.
\item The causal partial order: if $v \prec w$ in $G_0$, then $\Phi(v) \preceq \Phi(w)$ in $G_1$.
\item The BDG filter structure: the coarse-grained integers $(c_0^{(1)}, c_1^{(1)}, \ldots)$ inherit the signature and magnitudes of the underlying integers.
\end{itemize}

Applied iteratively, $G_2 = \Phi(G_1), G_3 = \Phi(G_2), G_4 = \Phi(G_3)$ define the four tiers. The same dynamical law operates at every level; only the effective degrees of freedom change.

\subsection{Markov blankets and decoherence}
\label{sec:markov_blanket}

A key structural fact from Paper~I: the interior of any causal subgraph is shielded from its exterior by its boundary. This shielding property (the Markov blanket theorem) is stated and Lean-verified at the discrete-graph level of Paper~I; the paper's Lean supplement contains the relevant theorem (\texttt{markov\_blanket}). \emph{Citation audit note:} the exact Lean file carrying this theorem is being reconfirmed against the audited Paper~I supplement following the file-attribution corrections applied to Paper~II (where \texttt{universe\_closure} and \texttt{confinement\_lengths} were found in \texttt{RA\_D1\_Proofs.lean} rather than \texttt{RA\_GraphCore.lean}). Readers should treat the theorem name as load-bearing; the file path is the current best attribution and will be corrected in the Lean supplement if needed. This shielding provides the structural mechanism for decoherence at every tier:

\begin{proposition}[Hierarchical decoherence]
\label{prop:hierarchical_decoherence}
At each coarse-graining level, a Tier-$k$ structure interacts with its environment only through its Tier-$(k-1)$ boundary. Interior dynamics are shielded from exterior influences except through the Markov blanket. The decoherence timescale at each tier follows from the actualization rate at the boundary, not from the total system size.
\end{proposition}

This explains why biological Tier-4 structures can maintain internal quantum coherence (e.g., in photosynthesis, olfactory receptors, or radical-pair magnetoreception) for timescales far exceeding what naive decoherence estimates would predict. The Markov blanket structurally isolates the functional subsystem from the thermal environment; decoherence acts only on the subset of degrees of freedom that cross the boundary.

The hierarchical decoherence picture gives RA a natural account of biological quantum effects: they are not mysterious exceptions to decoherence theory but structural consequences of the recursive coarse-graining of the causal graph.

%% ═══════════════════════════════════════════════════════════
\section{The Causal Firewall}
\label{sec:firewall}
%% ═══════════════════════════════════════════════════════════

\begin{definition}[Causal Firewall]
The Causal Firewall is the set of physical conditions under which a causal network can sustain recursive closure---self-maintaining boundary conditions through its own actualization dynamics. It is defined by two substrate-independent conditions.
\end{definition}

\textbf{Condition F1 (Quantum Darwinism redundancy).} The environment must record the outcomes of actualization events redundantly, so that the information content of each event is imprinted on macroscopic timescales rather than lost to quantum decoherence. Formally: the quantum Darwinism redundancy parameter $R_{\mathrm{QD}}\ge 1$, meaning each actualization event is independently witnessed by at least one environmental fragment.

\textbf{Condition F2 (Non-relativistic decoherence control).} The local thermal environment must be cold enough relative to the Unruh temperature $T_U=\hbar a/(2\pi c k_B)$ that the actualization rate is finite and controlled:
\begin{equation}
\frac{T_U}{T_{\mathrm{env}}} \ll 1.
\label{eq:f2}
\end{equation}
This is satisfied for all chemistry below nuclear energies. It excludes only extreme astrophysical environments: neutron star interiors, black hole vicinities, and the first microseconds after nucleation.

\begin{proposition}[Substrate independence]
Conditions F1 and F2 make no reference to molecular class, solvent chemistry, temperature range, or biological alphabet. Carbon-based life in aqueous solution is one realization. Any chemistry satisfying F1 and F2 can support Tier~4 complexity.
\end{proposition}

The physical content of the Causal Firewall is that it corresponds to the $\mu=1$ Erd\H{o}s--R\'enyi percolation threshold of the local causal graph---the onset of a giant connected component in which all observers agree on a shared causal ledger. Below this threshold, individual actualization events are causally isolated; above it, they form a self-reinforcing network capable of maintaining its own boundary conditions.

%% ═══════════════════════════════════════════════════════════
\section{The Universal Decoherence Timescale}
\label{sec:decoherence}
%% ═══════════════════════════════════════════════════════════

The Causal Firewall condition $\mu=\lambda\tau_d\ell^3\ge 1$ involves a decoherence timescale $\tau_d$ that can be derived from RA primitives at biological temperatures.

\subsection{The on-shell bath identification}

At temperature $T=300$\,K, two bath timescales are relevant:
\begin{itemize}[nosep]
\item \textbf{Vibrational (optical phonon) bath:} $\tau_c=\tau_{\mathrm{vib}}\approx 20$\,fs, $\hbar\omega_{\mathrm{vib}}/(k_BT)\approx 0.76$--$2.5$. These modes are on-shell ($\hbar\omega \gtrsim k_BT\Delta S^*$).
\item \textbf{Debye (acoustic) bath:} $\tau_c=\tau_{\mathrm{Debye}}\approx 8.3$\,ps, $\hbar\omega_{\mathrm{Debye}}/(k_BT)\approx 3\times 10^{-3}$. These modes are off-shell and do not contribute to actualization.
\end{itemize}

The $\Delta S^*$ criterion selects the vibrational bath: vibrational modes have $\hbar\omega$ comparable to $k_BT\Delta S^*\approx k_BT\times 0.6$, while acoustic modes have $\hbar\omega\ll k_BT\Delta S^*$. The filter identifies the correct bath, not an additional assumption.

\subsection{The universal decoherence timescale}

Applying the universal formula $\tau_d=\Delta S^*/(\Gamma_{\mathrm{opt}}\cdot\Delta S_{\mathrm{per\,phonon}})$ at $T=300$\,K:
\begin{equation}
\tau_d = \frac{\Delta S^*}{\Gamma_{\mathrm{opt}}\cdot(\hbar\omega_{\mathrm{opt}}/k_BT)} = \frac{0.601}{3.93\times 10^{13}\times 0.76} = 2.0\times 10^{-14}\;\mathrm{s} = 20\;\mathrm{fs}.
\end{equation}
This is the vibrational decoherence timescale. It is not a free parameter: $\Delta S^*=0.601$ is determined by the BDG integers, $\Gamma_{\mathrm{opt}}=k_BT/\hbar$ is determined by temperature, and $\Delta S_{\mathrm{per\,phonon}}$ is determined by the on-shell criterion.

\subsection{Non-Markovian corrections}

For slow-modulation environments (reorganization energy $\lambda_{\mathrm{reorg}}\ll k_BT$), the Kubo line-shape theory gives a non-Markovian correction:
\begin{equation}
\kappa = \frac{\sqrt{2\Delta S^*}}{\mu_{\mathrm{NM}}},
\end{equation}
where $\mu_{\mathrm{NM}}=\Delta\tau_c$ is the non-Markovian bath parameter. For RNA-world prebiotic chemistry ($\lambda_{\mathrm{reorg}}\approx 0.1$--$0.3$\,eV):
\begin{equation}
\kappa\approx 0.2\text{--}0.5, \qquad N_{\min}^{\mathrm{NM}}=\lceil 1/\kappa\rceil \approx 2\text{--}5.
\end{equation}

This is the range of minimum molecular chain lengths predicted by RA for the first self-replicating systems, consistent with the known minimum polymer lengths in RNA-world scenarios. The Markovian limit ($\kappa\to 1$) gives $N_{\min}=\lceil 1/\Delta S^*\rceil=2$; the non-Markovian correction raises this to 2--5 depending on the chemical environment.

%% ═══════════════════════════════════════════════════════════
\section{Assembly Depth and Assembly Theory}
\label{sec:assembly}
%% ═══════════════════════════════════════════════════════════

Assembly Theory~\cite{Sharma2023} measures molecular complexity via the assembly index~$A$: the minimum number of joining operations needed to construct a molecule from basic building blocks. RA defines an analogous quantity.

\begin{definition}[RA assembly depth]
The RA assembly depth $\tilde{A}_{\mathrm{RA}}$ is the minimum number of causally irreducible actualization layers needed to produce a molecular structure from elemental precursors.
\end{definition}

The relationship between the two indices is:
\begin{equation}
\tilde{A}_{\mathrm{RA}} \le A \le \tilde{A}_{\mathrm{RA}} \times b_{\max},
\end{equation}
where $b_{\max}$ is the maximum branching factor per layer. For molecules of moderate complexity, the two converge.

\textbf{Worked example: glycolysis.} The glycolysis pathway consists of 10 enzymatic reactions. Coarse-graining by causal irreducibility (grouping reactions whose products are used only by the next step) gives three causally irreducible layers: $\tilde{A}_{\mathrm{RA}}(\text{glycolysis})=3$. This sits just above the Causal Firewall minimum ($N_{\min}\ge 2$), consistent with glycolysis being among the oldest metabolic pathways.

\subsection{The E.~coli worked example}
\label{sec:ecoli}

The coarse-graining applies at every scale. We compute the RA assembly depth for \emph{E.~coli}, a well-characterized model organism.

The complete \emph{E.~coli} metabolic network (KEGG/BiGG) contains approximately $2{,}500$ enzymatic reactions across $\sim 40$ identifiable pathways. Coarse-graining at the pathway level (glycolysis, TCA cycle, pentose phosphate, fatty acid synthesis, amino acid biosynthesis, nucleotide biosynthesis, etc.) produces a depth hierarchy:

\begin{center}
\begin{tabular}{lcc}
\toprule
Coarse-graining level & Number of layers & $\tilde{A}_{\mathrm{RA}}$ contribution \\
\midrule
Individual reactions & $\sim 2{,}500$ & --- \\
Pathway-level & $\sim 40$ & 3--5 per pathway \\
Network-level (metabolic supergroups) & $\sim 8$ & 2--3 \\
Organism-level (the self-replicating cell) & 1 & 1 \\
\bottomrule
\end{tabular}
\end{center}

Applied recursively, the coarse-graining yields $\tilde{A}_{\mathrm{RA}}(\text{E. coli}) \approx 4$--$6$ layers, depending on the exact partition scheme. This places \emph{E.~coli} well above the origin-of-life minimum $N_{\min} \approx 2$--$5$, consistent with its status as a modern organism whose metabolic complexity has accumulated over $\sim 3.5 \times 10^9$ years of evolution.

The Assembly Theory index for an entire \emph{E.~coli} cell is not precisely known but is estimated to be $A \sim 10^{10}$--$10^{11}$ based on the complexity of its proteome and genome. The ratio $A / \tilde{A}_{\mathrm{RA}} \sim 10^9$ reflects the branching factor $b_{\max}$: each causally irreducible layer in a cell operates in parallel across billions of molecular degrees of freedom.

\subsection{Computational evidence from DFT calculations}
\label{sec:dft}

The F1 condition (quantum Darwinism redundancy) can be tested computationally via density functional theory (DFT) calculations on small prebiotic molecules. Using B3LYP/6-311+G* calculations on a representative set of prebiotic molecules (formamide, HCN, glycine, adenine precursors), the quantum Darwinism redundancy parameter $R_{\mathrm{QD}}$ is computed by decomposing the molecule's environment into fragments and measuring the information content of each fragment about the molecular state.

For all tested molecules, $R_{\mathrm{QD}} > 1$ at biological temperatures ($T = 300$\,K), indicating that F1 is satisfied. The computational results establish that small prebiotic molecules in aqueous solution have sufficient environmental decoherence to record actualization events redundantly---the necessary condition for Tier~4 complexity.

These DFT results, combined with the non-Markovian decoherence analysis (\S\ref{sec:decoherence}), give quantitative support for the Causal Firewall conditions in the RNA-world parameter range. The computational protocol is described in the accompanying supplementary material to RACI~\cite{Sandeman2026RACI}.

\paragraph{Audit status of the DFT/F1 artifact.} The B3LYP/6-311+G* protocol and the specific $R_{\mathrm{QD}}$ values for the representative molecule set are claimed at the CV (computation-verified) tier on the basis of the RACI supplementary computation. In the present Paper~IV suite audit, the underlying DFT output files were not re-run or independently inspected; the CV label rests on the supplementary material being reproducible from its stated protocol (B3LYP, 6-311+G*, representative molecule set, $T=300$\,K). Readers who require audit-level independence should regard this row as \emph{plausible and reproducible in principle but not audited within the present suite pass}, and should consult the RACI supplementary material for the underlying computation.

\begin{conjecture}[Assembly convergence]
For molecules with assembly index $A>15$, the RA assembly depth and the Assembly Theory index converge: $\tilde{A}_{\mathrm{RA}}\to A$ in the limit of low branching. This is testable using the KEGG/BiGG biochemical pathway database to compute both indices for metabolic networks.
\end{conjecture}

%% ═══════════════════════════════════════════════════════════
\section{The Origin-of-Life Sandwich Bound}
\label{sec:sandwich}
%% ═══════════════════════════════════════════════════════════

Any proposed prebiotic mechanism must produce a system with RA assembly depth within a computable range.

\textbf{Lower bound:} $\tilde{A}_{\mathrm{RA}}\ge 2$. An autocatalytic cycle requires at least two linked causally irreducible actualization layers forming a closed loop. A single layer cannot be self-replicating---it has no internal feedback.

\textbf{Upper bound:} $\tilde{A}_{\mathrm{RA}}\le A_{\max}(\ell, T, \lambda_{\mathrm{reorg}})$, where $A_{\max}$ is the maximum assembly depth sustainable at the local causal density, temperature, and bath coupling. This is computable from the non-Markovian decoherence timescale (\S\ref{sec:decoherence}) and the BDG causal graph topology at the Planck scale.

The sandwich bound constrains any proposed origin-of-life candidate:
\begin{equation}
2 \le \tilde{A}_{\mathrm{RA}}(M_0) \le A_{\max}(\ell, T, \lambda_{\mathrm{reorg}}),
\end{equation}
where $M_0$ is the first self-replicating molecule. RNA-world, metabolism-first, and crystal-template scenarios can each be quantitatively evaluated against this bound.

\textbf{Computational evidence.} B3LYP/6-311+G* density functional theory calculations supporting Condition~F1 have been performed for small prebiotic molecules, verifying that the quantum Darwinism redundancy parameter exceeds unity for molecular environments at biological temperatures~\cite{Sandeman2026RACI}.

%% ═══════════════════════════════════════════════════════════
\section{Quantum Information and Thermodynamics}
\label{sec:qinfo}
%% ═══════════════════════════════════════════════════════════

The RA framework gives specific predictions for quantum computing and resolves the foundational puzzles of quantum thermodynamics. Both topics follow from the same actualization-event primitive.

\paragraph{Scope note.} The results in this section---the Kinematic Coherence Bound, the spin-bath collapse timescale $t^*=0.274/g$, the Landauer bound, and the Maxwell's demon analysis---are inherited from the RAQI and RAQM companion papers rather than newly derived in Paper~IV. They are presented here for suite coherence, because the complexity-and-life architecture developed in \S\S\ref{sec:hierarchy}--\ref{sec:sandwich} uses the KCB-style quantum-information machinery as its decoherence backbone. Readers should not read the presentation in this section as fresh first-principles closure within Paper~IV; the load-bearing derivations live in RAQI/RAQM.

\subsection{The Kinematic Coherence Bound}
\label{sec:kcb}

A key prediction for quantum computing: fault-tolerant quantum arrays have a hard structural ceiling on their size, set by the accumulated actualization rate across all physical qubits.

\begin{proposition}[Kinematic Coherence Bound (KCB)]
\label{prop:kcb}
For a quantum array of $N$ physical qubits with single-qubit error probability $p$, fault-tolerant operation requires:
\begin{equation}
N \leq N_{\max} = \frac{\eta}{p_{\mathrm{th}}},
\end{equation}
where $p_{\mathrm{th}} = \alpha_{\mathrm{EM}} \approx 1/137$ is the RA-predicted error threshold and $\eta$ is the single-qubit quality factor.
\end{proposition}

The logic: each qubit is an additional kinematic site at which actualization events can occur. The RA minimum BDG stability for electrons and photons (standard qubit substrates) is exactly 1 (Lean-verified, L12 in Paper~II). The cumulative probability of spurious actualizations across $N$ qubits exceeds the fault-tolerance budget when $N \cdot p > \eta \cdot p_{\mathrm{th}}$.

For current superconducting qubit technologies ($\eta \sim 10$--$100$, $p_{\mathrm{th}} \sim 10^{-2}$), this predicts hard ceilings at $N_{\max} \sim 10^3$--$10^4$ logical qubits. Beyond this size, increasing fidelity does not rescue fault tolerance---the bound is kinematic, not just operational.

\textbf{Distinctive prediction:} Standard decoherence theory predicts exponentially difficult but unbounded scaling. RA predicts a \emph{hard ceiling} at a specific size determined by hardware quality factors. Near-future quantum error correction experiments at the 1000-logical-qubit scale should be sensitive to this distinction.

\subsection{Decoherence threshold behavior}
\label{sec:threshold_decoherence}

RA predicts qualitatively different decoherence below and above the on-shell threshold $\Delta S^*$. Below the threshold, environmental interactions do not actualize and therefore do not contribute to decoherence. Above it, environmental interactions actualize and cause decoherence.

\begin{conjecture}[Virtual/real distinction for QEC]
Operating qubits in a regime where the qubit--environment coupling energy is below the on-shell threshold $\Delta S^* \cdot k_B T \approx 0.601 k_B T$ should produce qualitatively suppressed decoherence---not just quantitatively reduced, but mechanistically different from standard Lindblad predictions.
\end{conjecture}

Superconducting qubits operating at millikelvin temperatures are already in this regime; the RA prediction is that their decoherence rates should show threshold behavior absent in standard Lindblad master equation predictions. Direct tests would vary the coupling strength through the threshold and observe the transition.

\subsection{Spin-bath collapse timescale}
\label{sec:spinbath}

For a quantum system coupled to a spin bath with externally measured coupling strength $g$, RA predicts a specific timescale for wave function collapse:
\begin{equation}
\label{eq:spinbath}
\boxed{t^* = \frac{1}{g}\sqrt{\frac{\Delta S^*}{2}} \approx \frac{0.274}{g}.}
\end{equation}
The numerical prefactor carries no additional RA free parameters: $\Delta S^* = 0.601$ nats is determined by the BDG integers, and the $\sqrt{\Delta S^*/2}$ factor is fixed by the spin-bath dynamical model.

This distinguishes RA from GRW (Ghirardi--Rimini--Weber) and CSL (continuous spontaneous localization) collapse models, which predict observable collapse timescales with different parameter dependences. Superconducting circuit experiments can probe this regime by controlling the effective bath coupling $g$.

\subsection{Landauer's principle from the causal DAG}
\label{sec:landauer}

Landauer's principle states that erasing one bit of information requires dissipating at least $k_B T \ln 2$ of heat. In RA, this follows directly from the causal graph structure.

An information-bearing state corresponds to a set of actualized vertices in the causal graph. Erasing a bit means overwriting that state with a reference state, which requires actualizing new vertices that replace the old ones. The minimum number of actualization events required to ``overwrite'' one bit of information is one (the actualization of a single vertex carries $\log 2$ nats of information, the minimum non-trivial information content).

The thermodynamic cost of one actualization event at temperature $T$ is $k_B T \Delta S^* \approx 0.6 k_B T$. For a process that overwrites information rather than creating it, the cost must be at least $k_B T \ln 2 \approx 0.693 k_B T$---the minimum needed to decouple the overwritten vertex from the original information. Landauer's principle is the thermodynamic floor imposed by the actualization threshold.

\begin{proposition}[RA derivation of Landauer bound]
\label{prop:landauer}
The Landauer bound $\Delta E \geq k_B T \ln 2$ per bit erased is the thermodynamic cost of actualization at the minimum non-trivial information content. It follows from the BDG actualization threshold $\Delta S^*$ applied at the bit-erasure scale.
\end{proposition}

\subsection{Maxwell's demon dissolved}
\label{sec:maxwell}

Maxwell's demon was a thought experiment designed to show that the second law of thermodynamics could be violated by an intelligent being sorting gas molecules by velocity. The standard resolution attributes the second law's survival to the information-processing cost the demon incurs.

In RA, the demon's information-processing cost is not a separate postulate: it is the actualization cost of the demon's observations and decisions. Each measurement the demon makes about a molecule is an actualization event with thermodynamic cost $\geq k_B T \Delta S^*$. Each decision-bit the demon erases (to free up memory) incurs the Landauer cost $k_B T \ln 2$.

The demon's total thermodynamic cost exceeds the entropy reduction he achieves by sorting molecules, as required by the second law. The apparent paradox dissolves because information processing is not free---it is the physical process of writing actualization events into the causal graph, with associated thermodynamic cost.

\subsection{The arrow of time as structural}
\label{sec:qinfo_arrow}

In standard thermodynamics, the arrow of time is statistical: it emerges from the tendency of closed systems to increase entropy. In RA, as developed in Paper~I, the arrow of time is structural---it is the growth direction of the causal DAG, which is primitive.

The connection to quantum information: an information-bearing state is a pattern of actualized vertices. ``Forgetting'' a state means actualizing new vertices that overwrite the old ones, but the old vertices are not removed---they are simply no longer referenced by the current state. The causal graph is append-only; history is preserved in principle even when it is inaccessible in practice. This is the structural origin of the thermodynamic arrow: information can be destroyed in an observer's frame (by being decoupled from current dynamics) but never erased from the global graph.

%% ═══════════════════════════════════════════════════════════
\section{Biosignature Criteria and SETI Predictions}
\label{sec:seti}
%% ═══════════════════════════════════════════════════════════

The Causal Firewall gives three universal biosignature criteria that do not presuppose any specific chemistry:

\textbf{B1 (Redundant information recording):} The environment satisfies F1---it supports quantum Darwinism, meaning actualized states are redundantly recorded in environmental degrees of freedom.

\textbf{B2 (Controlled decoherence):} The environment satisfies F2---the Unruh temperature is far below the ambient thermal temperature, ensuring finite and controlled actualization rates.

\textbf{B3 (Edge-of-chaos density):} The local causal actualization density operates near $\mu\approx 1$---the percolation threshold at which the causal graph transitions from isolated clusters to a giant connected component.

\subsection{Falsifiable predictions}

\textbf{Structural impossibility:} Life is structurally impossible in environments where $T_U\sim T_{\mathrm{env}}$. This excludes neutron star surfaces, black hole vicinities, and the first microseconds after cosmological nucleation. The prediction is categorical, not statistical.

\textbf{Substrate universality:} All extraterrestrial life---regardless of molecular substrate---satisfies F1, F2, and B3. This predicts that SETI searches weighted by B1--B3 will outperform searches weighted by conventional habitable-zone criteria (liquid water, Earth-like temperature).

\textbf{Causal depth as selection pressure:} Causal depth is an independent selection pressure in natural selection, distinct from energy efficiency and reproductive rate. Organisms with greater causal depth (more actualization layers between environmental input and behavioral output) have a structural advantage in complex environments. This is testable via comparative analysis of causal depth across taxa.

\textbf{Assembly convergence:} The RA assembly depth and the Assembly Theory index converge for molecules with $A>15$. Systematic computation across the KEGG/BiGG metabolic network would test this conjecture.

%% ═══════════════════════════════════════════════════════════
\section{Perception, Agency, and Consciousness: Structural Constraints}
\label{sec:consciousness}
%% ═══════════════════════════════════════════════════════════

\textbf{Epistemic flag.} This is the most exploratory section of the most exploratory paper in the suite. RA does not claim to solve the hard problem of consciousness. What follows is a structural framework---a set of necessary conditions and exclusions that RA's causal architecture imposes---not a theory of phenomenal experience. The reader should treat everything in this section as argued (AR) or conjectural (CN) unless explicitly stated otherwise.

The complexity hierarchy of \S\ref{sec:hierarchy} extends naturally to the question of perception, agency, and consciousness. We organize the discussion around four categories: structural consequences of the framework's causal architecture, structural exclusions, structural permissions, and conjectures.

\subsection{Structural consequence: perception as actualization anchoring}

A perceiving system does not passively receive information; it \emph{actualizes} sensory states. When a photon enters an eye and triggers a photoreceptor cascade, the cascade writes a sequence of actualization events into the observer's causal graph. The percept is not the external stimulus but the actualization pattern in the observer's Tier-4 network.

Three structural consequences follow from this identification (all DR within RA):
\begin{itemize}[nosep]
\item \textbf{Perception is first-person in principle.} Each observer's perceptual graph is causally distinct from every other observer's, because actualization events are local to the perceiving subgraph.
\item \textbf{Perception is irreversible.} Once a percept is written into the causal graph, it cannot be undone. This is the RA-native origin of the phenomenological arrow of experience.
\item \textbf{Perception is rate-limited.} The actualization frequency of a perceiving subsystem is bounded by its BDG boundary; no Tier-4 system can perceive faster than its Markov-blanket actualization rate permits.
\end{itemize}

\subsection{Structural consequence: agency as recursive causal closure}

An agent is a Tier-4 pattern that maintains its own boundary conditions through its own actualization dynamics---a pattern whose future states depend on its own past states, via actualization events that the pattern itself generates rather than receiving passively from the environment. Agency in this sense is a matter of degree, not kind: all Tier-4 systems exhibit some recursive closure; the degree of agency depends on the proportion of internally generated actualization events. This is a structural consequence of the tier hierarchy (DR), not an additional postulate.

\subsection{What RA excludes: panpsychism}

RA's causal architecture excludes panpsychist frameworks that attribute phenomenal character to simple physical systems.

The argument is structural: a single actualization event, or a small cluster of them, does not have recursive causal closure. Consciousness---whatever it turns out to be---requires a Tier-4 system whose internal dynamics are rich enough to sustain self-referential processing through multiple coarse-graining layers. A rock, a thermostat, or a single photon does not meet this criterion, because none of them maintains its own boundary conditions through its own actualization activity. This distinguishes RA from Integrated Information Theory (IIT)~\cite{Tononi2016} and constitutive panpsychism, both of which attribute phenomenal character to systems that RA classifies as Tier-2 or Tier-3 structures lacking recursive closure.

\textbf{Status:} This exclusion is a structural consequence (DR) of the tier hierarchy and the Causal Firewall conditions, not an independent postulate about consciousness.

\subsection{What RA permits: artificial consciousness}

RA's conditions for consciousness are substrate-independent: F1 (redundant information recording), F2 (controlled decoherence), and recursive causal closure above a sufficient depth threshold. Any physical substrate that implements a Tier-4 network satisfying these conditions can, in principle, support whatever causal architecture consciousness requires---whether the substrate is biological, silicon, or something else entirely.

\textbf{Status:} This is a structural permission (DR), not a prediction. RA does not predict that any existing artificial system is conscious, nor does it specify the depth threshold above which consciousness arises.

\subsection{What RA conjectures: the depth threshold}

\begin{conjecture}[Structural necessary condition for consciousness]
\label{conj:consciousness}
A causal network exhibits phenomenal consciousness only if: (i) it satisfies F1 and F2 (Causal Firewall), and (ii) its RA assembly depth exceeds a threshold $\tilde{A}_{\mathrm{RA}}^* \gtrsim 5$.
\end{conjecture}

The threshold value $\tilde{A}_{\mathrm{RA}}^* \approx 5$ is estimated from the observed minimum complexity of organisms exhibiting neurally encoded perception (roughly the complexity level of simple arthropods with centralized nervous systems). It is \emph{not} derived from the BDG integers or from any first-principles argument within RA. It is a conjecture motivated by comparative neuroscience, not a consequence of the framework.

Even if the threshold is correct, it provides only a \emph{necessary} condition, not a sufficient one. Many systems may exceed $\tilde{A}_{\mathrm{RA}}^* = 5$ without being phenomenally conscious. The hard problem of consciousness---why there is something it is like to be a conscious system---remains open. RA narrows the structural space in which the answer must lie; it does not supply the answer.

%% ═══════════════════════════════════════════════════════════
\section{Epistemic Status Summary}
\label{sec:status}
%% ═══════════════════════════════════════════════════════════

\begin{center}
\begin{tabular}{llc}
\toprule
Result & Basis & Status \\
\midrule
\multicolumn{3}{l}{\textit{Complexity hierarchy}} \\
Four-tier complexity hierarchy & Recursive coarse-graining $\Phi$ & DR \\
Hierarchical decoherence & Markov blanket theorem (LV) & DR \\
$\mu=1$ as unifying architecture across 5 contexts & BDG Poisson-CSG + Erd\H{o}s--R\'enyi & Mixed (see below) \\
\quad — Context 1 (QCD confinement) & BDG path-weight enumeration (Paper~II) & CV \\
\quad — Context 2 (galactic rotation/dim. reduction) & Paper~III continuum-bridge (DJW theorem) & DR conditional$^\dagger$ \\
\quad — Context 3 (QC error threshold $p_{\mathrm{th}}$) & KCB identification & PI \\
\quad — Context 4 (actualization threshold $\Delta S^*$) & Paper~I D4U02 enumeration & CV \\
\quad — Context 5 (Causal Firewall percolation) & This paper, F1+F2+Erd\H{o}s--R\'enyi & DR (this paper) \\
\midrule
\multicolumn{3}{l}{\textit{Causal Firewall}} \\
Causal Firewall (F1 + F2) & Substrate-independent conditions & DR \\
$\tau_d=20$\,fs at $T=300$\,K & $\Delta S^*$ + on-shell bath & CV \\
$\kappa\approx 0.2$--$0.5$ (RNA-world) & Kubo line-shape theory + $\lambda_{\mathrm{reorg}}$ & CV \\
$N_{\min}=2$--$5$ & $\lceil 1/\kappa\rceil$ & CV \\
Sandwich bound on first replicator & F1 + F2 + $N_{\min}$ & DR \\
B3LYP support for F1 & DFT computation (RACI supplement) & CV (not re-audited in suite pass) \\
\midrule
\multicolumn{3}{l}{\textit{Assembly theory connection}} \\
$\tilde{A}_{\mathrm{RA}}(\text{glycolysis})=3$ & Pathway coarse-graining & CV \\
$\tilde{A}_{\mathrm{RA}}(\text{E. coli})\approx 4$--$6$ & Metabolic network coarse-graining & CV \\
Assembly convergence for $A>15$ & Structural conjecture & Conjecture \\
\midrule
\multicolumn{3}{l}{\textit{Quantum information}} \\
KCB: $N_{\max}=\eta/p_{\mathrm{th}}$ & Cumulative actualization rate & DR \\
Threshold decoherence behavior & Virtual/real distinction & Conjecture \\
$t^*=0.274/g$ (spin-bath collapse) & Parameter-free from $\Delta S^*$ & DR \\
Landauer bound $\geq k_B T \ln 2$ & Actualization thermodynamics & DR \\
Maxwell's demon dissolved & Information = actualization cost & DR \\
Arrow of time structural, not statistical & Graph growth direction & DR \\
\midrule
\multicolumn{3}{l}{\textit{Biosignatures and SETI}} \\
Biosignature criteria B1--B3 & F1 + F2 + $\mu=1$ & DR \\
Life impossible at $T_U\sim T_{\mathrm{env}}$ & F2 structural exclusion & Prediction \\
Substrate universality of life & F1 + F2 independence & Prediction \\
Causal depth as selection pressure & Tier~4 recursive closure & AR \\
\midrule
\multicolumn{3}{l}{\textit{Perception and consciousness}} \\
Perception as actualization anchoring & Structural consequence of Tier-4 & DR \\
Agency as recursive causal closure & Structural consequence of Tier-4 & DR \\
Panpsychism excluded & Recursive closure required (structural) & DR \\
Artificial consciousness permitted & Substrate-independent conditions & DR \\
Depth threshold $\tilde{A}_{\mathrm{RA}}^* \gtrsim 5$ & Comparative neuroscience estimate & CN \\
\bottomrule
\end{tabular}
\end{center}

\textbf{Legend:} LV = Lean-verified theorem. DR = derived (all steps explicit within RA). CV = computation-verified (exact derivation plus numerical evaluation). PI = phenomenological interpolation (motivated identification supported by numerical evidence, not derived from first principles). AR = argued (physically motivated, not all steps explicit). CN = conjecture (stated path to closure, no current derivation). Prediction = falsifiable consequence testable with current or near-future data.

%% ═══════════════════════════════════════════════════════════
\section{Conclusion}
\label{sec:conclusion}
%% ═══════════════════════════════════════════════════════════

This paper has extended the framework of Relational Actualism into the domains of complexity, life, agency, and cognition. Its central claim is not that these domains have already been reduced to a finished theorem set, but that they can now be discussed within a common RA-native ontology: finite actual graphs, structured potentia, shielding, recursive closure, and the causal firewall. In that sense, the paper's main contribution is architectural. It shows that the same discrete principles used elsewhere in the suite to treat matter, gravity, and severance can also be brought to bear on questions of stable complexity, biological organization, and higher-order cognitive structure.

The strongest results of the paper are structural. Recursive closure, shielding, and boundary formation emerge as the key organizational conditions under which stable complexity can persist. The causal firewall then provides a common language for discussing why many systems fail to cross from passive organization into active, self-maintaining, information-bearing complexity. In this respect, the paper's most durable claims concern organizational constraints, not full biological or cognitive closure. The life programme developed here should therefore be read as a proposal for the conditions under which life-like organization becomes possible within RA, not as a finished derivation of all biological structure.

The same caution applies to the discussions of agency, perception, and consciousness. The paper does not claim a sufficient-condition theory of phenomenal consciousness. Rather, it identifies a set of structural features that appear to be necessary, or at least strongly indicated, within the RA framework: shielding, recursive closure, persistent boundary organization, and stable actualization--perception loops. Conversely, the framework appears to disfavor pictures in which consciousness is attributed independently of such organization. These are significant constraints, but they do not yet amount to full closure on the relation between complexity, cognition, and experience.

The empirical and observational claims of the paper should be read in the same graded way. Some are structural forecasts of the framework; some are concrete but still exploratory predictions concerning biosignatures, complexity thresholds, and the causal firewall; and some remain conjectural extensions that will require substantial further work before they can be regarded as robust.

What the paper establishes, then, is not a complete theory of life or mind, but a disciplined RA-native programme for approaching them. It argues that complexity, agency, and cognition need not be appended to the earlier papers as alien domains; they can be treated as higher-order organizational consequences of the same discrete architecture. That is already a meaningful result.

Several major open problems remain. On the life side, the proposed thresholds for stable biological organization require sharper quantitative development. On the agency side, the relation between recursive closure and robust action remains partly schematic. On the consciousness side, the most important missing step is clear: the framework has not yet produced a sufficient-condition criterion, and the conjectural thresholds discussed here remain to be justified or replaced. These are not defects to be hidden; they are the next substantive tests of whether the RA complexity programme can advance from a suggestive architecture to a predictive theory.

What RA offers at present is not closure, but a disciplined narrowing of the problem.

%% ═══════════════════════════════════════════════════════════
\section*{Acknowledgments}
%% ═══════════════════════════════════════════════════════════

This work was developed in collaboration with Claude (Anthropic, Opus 4.6) for computation, decoherence timescale derivations, and manuscript preparation, and with ChatGPT (OpenAI, GPT-4o) for the non-Markovian bath analysis and structural review. The ontological commitments, physical reasoning, and framework direction are the author's.

\begin{thebibliography}{99}
\bibitem{Benincasa2010} D.~M.~T.~Benincasa and F.~Dowker, ``The scalar curvature of a causal set,'' \emph{Phys. Rev. Lett.} \textbf{104}, 181301 (2010).
\bibitem{Erdos1960} P.~Erd\H{o}s and A.~R\'enyi, ``On the evolution of random graphs,'' \emph{Publ. Math. Inst. Hung. Acad. Sci.} \textbf{5}, 17 (1960).
\bibitem{Sharma2023} A.~Sharma et al., ``Assembly theory explains and quantifies selection and evolution,'' \emph{Nature} \textbf{622}, 321 (2023).
\bibitem{Zurek2009} W.~H.~Zurek, ``Quantum Darwinism,'' \emph{Nat. Phys.} \textbf{5}, 181 (2009).
\bibitem{Kubo1969} R.~Kubo, ``A stochastic theory of line shape,'' \emph{Adv. Chem. Phys.} \textbf{15}, 101 (1969).
\bibitem{Sandeman2026RACI} J.~F.~Sandeman, ``Relational Actualism: Complexity and Intelligence (RACI),'' Preprint, 2026. doi.org/10.5281/zenodo.19198224.
\bibitem{Sandeman2026RAHC} J.~F.~Sandeman, ``Algorithmic Causal Dynamics (RAHC),'' Preprint, 2026. doi.org/10.5281/zenodo.19198171.
\bibitem{Sandeman2026RACF} J.~F.~Sandeman, ``The Causal Firewall (RACF),'' \emph{Int. J. Astrobiol.} (submitted), 2026. doi.org/10.5281/zenodo.19319374.
\bibitem{Landauer1961} R.~Landauer, ``Irreversibility and heat generation in the computing process,'' \emph{IBM J. Res. Dev.} \textbf{5}, 183 (1961).
\bibitem{Bennett1982} C.~H.~Bennett, ``The thermodynamics of computation---a review,'' \emph{Int. J. Theor. Phys.} \textbf{21}, 905 (1982).
\bibitem{Tononi2016} G.~Tononi, M.~Boly, M.~Massimini, and C.~Koch, ``Integrated information theory: from consciousness to its physical substrate,'' \emph{Nat. Rev. Neurosci.} \textbf{17}, 450 (2016).
\bibitem{Sandeman2026RAQI} J.~F.~Sandeman, ``Relational Actualism: Implications for Quantum Information (RAQI),'' Preprint, 2026. doi.org/10.5281/zenodo.19198285.
\end{thebibliography}

\end{document}
